% ID: SEL-D0073-OB-E01
% ESTADO: cotejada. Segunda pasada visual y borrador compilado revisados.
% Fuente: Selectividad/sel_2008_matematicas_aplicadas/Mat CCSS 1 - Examen y criterios de correción.pdf
% SHA-256: 2461e302e686c60322df8edd0a004b11753e95c4824dc5835b16b76c05174e81
% Página: 2; opción B; ejercicio original 1.

\textbf{(3 puntos)} Un pastelero dispone de 150 kg de harina, 22 kg de azúcar y 26 kg de mantequilla para hacer dos tipos de tartas, A y B. Para hacer una hornada de tartas del tipo A se necesitan 3 kg de harina, 1 kg de azúcar y 1 kg de mantequilla, mientras que para hacer una hornada de tartas del tipo B se necesitan 6 kg de harina, $0.5$ kg de azúcar y 1 kg de mantequilla. Sabiendo que el beneficio que se obtiene al vender una hornada del tipo A es de 20 \texteuro{} y de 30 \texteuro{} al vender una hornada del tipo B, determine cuántas hornadas de cada tipo debe hacer y vender para maximizar sus beneficios.
